\chapter{Yangians and Shifted Yangians}
\label{chap:yangians}

This chapter treats Yangians and their shifted variants, with connections to 
Coulomb branches and cohomological Hall algebras.


\section{Yangian Vertex Algebra}
\label{sec:yangian-vertex}

\subsection{Yangian Definition}

\begin{definition}[Yangian]\label{def:yangian}
The \textbf{Yangian} $Y(\mathfrak{g})$ associated to a simple Lie algebra $\mathfrak{g}$ 
is the associative algebra generated by $J^{(r)}_a$ for $a = 1, \ldots, \dim\mathfrak{g}$ 
and $r \geq 0$, with relations:
\begin{align*}
[J^{(0)}_a, J^{(0)}_b] &= f_{ab}^c J^{(0)}_c \\
[J^{(0)}_a, J^{(r)}_b] &= f_{ab}^c J^{(r)}_c \\
[J^{(1)}_a, [J^{(1)}_b, J^{(0)}_c]] - [J^{(0)}_a, [J^{(1)}_b, J^{(1)}_c]] &= \hbar^2 A_{abc}^{def} \{J^{(0)}_d, J^{(0)}_e, J^{(0)}_f\}
\end{align*}
where $\{x,y,z\} = \frac{1}{6}\sum_{\sigma \in S_3} x_{\sigma(1)}y_{\sigma(2)}z_{\sigma(3)}$ is the symmetrized triple product ($|S_3| = 6$), and $A_{abc}^{def}$ are structure constants determined by $\mathfrak{g}$ (Drinfeld's ``terrific'' relation).
\end{definition}

\begin{definition}[Yangian Current]\label{def:yangian-current}
The \textbf{Yangian current} is the generating function:
\[
J_a(u) = \sum_{r \geq 0} J^{(r)}_a u^{-r-1}
\]
\end{definition}

\begin{definition}[RTT Presentation]\label{def:yangian-rtt}
The \textbf{RTT presentation} (Yang--Baxter presentation) of
$Y(\mathfrak{g})$ is generated by a \emph{transfer matrix}
$T(u) = (T_{ij}(u))_{i,j=1}^{\dim V}
\in \mathrm{End}(V) \otimes Y(\mathfrak{g})$
(valued in $\mathrm{End}(V)$ for a fundamental representation~$V$,
not in~$\mathfrak{g}$), satisfying the \textbf{RTT relation}:
\begin{equation}\label{eq:rtt-relation}
R_{12}(u-v)\, T_1(u)\, T_2(v)
\;=\; T_2(v)\, T_1(u)\, R_{12}(u-v)
\end{equation}
where $R_{12}(u) = 1 - \hbar P_{12}/u + O(u^{-2})$ is the
Yang $R$-matrix, $P_{12}$ is the permutation, and the subscripts
$1, 2$ indicate the tensor factor in
$\mathrm{End}(V) \otimes \mathrm{End}(V) \otimes Y(\mathfrak{g})$.
Expanding in modes
$T_{ij}(u) = \delta_{ij} + \hbar \sum_{r \geq 0} T_{ij}^{(r)} u^{-r-1}$,
the RTT relation \eqref{eq:rtt-relation} is equivalent to:
\begin{equation}\label{eq:rtt-modes}
[T_{ij}^{(r+1)}, T_{kl}^{(s)}]
- [T_{ij}^{(r)}, T_{kl}^{(s+1)}]
\;=\; T_{kj}^{(r)} T_{il}^{(s)}
- T_{kj}^{(s)} T_{il}^{(r)}.
\end{equation}
The RTT relation is \emph{quadratic} in the entries of $T$, making
this presentation directly amenable to quadratic Koszul duality
(mechanism~(iii) of Remark~\ref{rem:three-koszul-mechanisms}).
The isomorphism with the Drinfeld presentation
(Definition~\ref{def:yangian}) is given by
$J_a^{(0)} = \sum_{i,j} \rho(e_a)_{ij}\, T_{ij}^{(0)}$
and higher modes from the coproduct structure.
\end{definition}


\subsection{Yangian Chiral Algebra}

\begin{construction}[Yangian Vertex Algebra]\label{constr:yangian-va}
The \textbf{Yangian chiral algebra} $Y(\mathfrak{g})^{\mathrm{ch}}$ is constructed as follows:

\textbf{State space:} Completion of $Y(\mathfrak{g})$-modules.

\textbf{Vertex operators:} $Y(J^{(r)}_a, z) = \sum_n J^{(r)}_{a,n} z^{-n-r-1}$.

\textbf{OPE:} The RTT relation translates to:
\[
J_a(z) J_b(w) = \frac{f_{ab}^c J_c(w)}{z - w} + \frac{\alpha_{ab}^{cd} {:}J_c(w) J_d(w){:}}{(z-w)^2} + \text{regular}
\]
with higher poles from the quadratic Serre relations.
\end{construction}

\begin{theorem}[Yangian as $\Eone$-Chiral; \ClaimStatusProvedHere]\label{thm:yangian-e1}
In the RTT presentation, the Yangian $Y(\mathfrak{g})^{\mathrm{ch}}$ carries the structure of a quadratic $\Eone$-chiral algebra (Definition~\ref{def:e1-chiral}).  The $\Eone$-locality axiom is $R$-twisted commutativity:
\begin{equation}\label{eq:r-twisted-locality}
J_a(z) J_b(w) - R^{ab}_{cd}(z-w)\, J_d(w) J_c(z) = 0.
\end{equation}
For nontrivial $R \neq P$ (the permutation), this is strictly $\Eone$,
not $\Einf$.
\end{theorem}

\begin{proof}
\textbf{Step 1: Quadraticity.}  The RTT relation~\eqref{eq:rtt-relation} is manifestly quadratic in the entries $T_{ij}(u)$.  More precisely, expanding $T_1(u) T_2(v)$ and $T_2(v) T_1(u)$ in the mode expansion $T_{ij}(u) = \delta_{ij} + \hbar \sum_r T_{ij}^{(r)} u^{-r-1}$, the relation~\eqref{eq:rtt-modes} expresses a quadratic relation among the generators $\{T_{ij}^{(r)}\}$.

\textbf{Step 2: $\Eone$-chiral structure.}
The vertex operators $Y(T_{ij}^{(r)}, z) = \sum_n T_{ij,n}^{(r)} z^{-n-r-1}$ define fields on a curve $X$.  The OPE:
\[
T_{ij}(z,u)\, T_{kl}(w,v) \sim \frac{R^{ij,kl}(u-v)}{z-w} \cdot (\text{fusion})
\]
is $R$-matrix-valued, with simple poles.  The $R$-matrix $R(u) = 1 + \hbar r/u + O(u^{-2})$ provides the braiding: $\sigma_{12} \circ \mu = R_{12} \cdot \mu \circ \sigma_{12}$, which is the $\Eone$-locality axiom.  Associativity follows from the Yang--Baxter equation $R_{12}R_{13}R_{23} = R_{23}R_{13}R_{12}$.

\textbf{Step 3: Non-commutativity.}
For $\mathfrak{g} = \mathfrak{sl}_2$, the classical $r$-matrix $r = \frac{1}{2}(e \otimes f + f \otimes e) + \frac{1}{4} h \otimes h$ (the split Casimir) satisfies $r \neq 0$, so the $R$-matrix $R(u) = 1 + \hbar r/u + O(u^{-2})$ is non-trivial: $R(u) \neq P$ for finite~$u$, and the RTT relation does not reduce to commutativity.  Thus the algebra is strictly $\Eone$.
\end{proof}

\begin{remark}[Presentations and quadraticity]\label{rem:yangian-presentations}
Whether the Yangian is ``quadratic'' depends on the presentation:
\begin{itemize}
\item In the \textbf{RTT presentation} (Definition~\ref{def:yangian-rtt}), the
defining relation \eqref{eq:rtt-relation} is
quadratic in $T$, so standard quadratic Koszul duality (mechanism~(iii) of
Remark~\ref{rem:three-koszul-mechanisms}) applies directly.
\item In the \textbf{Drinfeld presentation} (Definition~\ref{def:yangian}), the
``terrific relation'' is cubic in the generators, so Koszul duality requires
nilpotent completion (Appendix~\ref{app:nilpotent-completion}).
\end{itemize}
The two presentations yield isomorphic algebras but different Koszul dual
computations.  A rigorous construction of the $\Eone$-chiral structure and its
Koszul dual should specify which presentation is used.
\end{remark}


\section{Koszul Dual of the Yangian}\label{sec:yangian-koszul}

\subsection{RTT Bar Complex Computation}

\begin{theorem}[Yangian Bar Complex via RTT; \ClaimStatusProvedHere]\label{thm:yangian-bar-rtt}
For $\mathfrak{g} = \mathfrak{sl}_2$, in the RTT presentation, the bar complex $\bar{B}(Y(\mathfrak{sl}_2)^{\mathrm{ch}})$ is computed as follows.

\textbf{Generators:} In bar degree~1, the generators are the duals $T_{ij}^{(r)*}$ of the RTT generators, placed on $\overline{C}_2(X)$ with log forms.

\textbf{Differential:} The bar differential extracts the RTT relation via double residues:
\begin{equation}\label{eq:yangian-bar-diff}
d(T_{ij}^{(r)*} \otimes \eta_{12}) = \sum_{k} \bigl(T_{ik}^{(r)*} \otimes T_{kj}^{(0)*} - T_{ik}^{(0)*} \otimes T_{kj}^{(r)*}\bigr) \otimes \eta_{12} \wedge \eta_{23}
\end{equation}
This is the quadratic dual of the RTT relation~\eqref{eq:rtt-modes}: the bar differential encodes the orthogonal complement of the relations under the chiral pairing.

\textbf{Bar complex structure:} Since the RTT relation is quadratic, the Koszul dual $A^!$ is determined by the data in bar degrees $\leq 2$ alone: generators $V^*$ and relations $R^\perp$.  If the algebra is Koszul (i.e., $\mathrm{Ext}^{i,j}_A(\mathbb{C}, \mathbb{C}) = 0$ for $i \neq j$), then the bar-cobar quasi-isomorphism recovers $A$ from $A^!$.
\end{theorem}

\begin{proof}
The RTT relation~\eqref{eq:rtt-modes} is:
$[T_{ij}^{(r+1)}, T_{kl}^{(s)}] - [T_{ij}^{(r)}, T_{kl}^{(s+1)}] = T_{kj}^{(r)} T_{il}^{(s)} - T_{kj}^{(s)} T_{il}^{(r)}$.
This is quadratic in $T$, so the bar complex of the quadratic algebra has:
\begin{itemize}
\item $\bar{B}^0 = \mathbb{C}$ (vacuum),
\item $\bar{B}^1 = V^* = \mathrm{span}\{T_{ij}^{(r)*}\}$ (dual generators on $\overline{C}_2(X)$),
\item $\bar{B}^2 = R^\perp \subset V^* \otimes V^*$ (orthogonal complement of relations on $\overline{C}_3(X)$).
\end{itemize}
The differential $d\colon \bar{B}^2 \to \bar{B}^1$ is the dual of the inclusion $R \hookrightarrow V \otimes V$, which gives~\eqref{eq:yangian-bar-diff}.
\end{proof}

\subsection{Koszul Dual Identification}

\begin{theorem}[Yangian Koszul Dual; \ClaimStatusProvedHere]\label{thm:yangian-koszul-dual}
The Koszul dual of the Yangian $Y(\mathfrak{sl}_2)^{\mathrm{ch}}$ in the RTT presentation is the Yangian with inverse $R$-matrix:
\begin{equation}\label{eq:yangian-koszul-dual}
Y(\mathfrak{sl}_2)^{!,\mathrm{ch}} \;\cong\; Y_{R^{-1}}(\mathfrak{sl}_2)^{\mathrm{ch}}
\end{equation}
where $Y_{R^{-1}}$ is the RTT algebra defined by the relation:
\begin{equation}
R^{-1}_{12}(u-v)\, T_1^!(u)\, T_2^!(v) = T_2^!(v)\, T_1^!(u)\, R^{-1}_{12}(u-v)
\end{equation}
with $R^{-1}(u) = 1 + \hbar P/u + O(u^{-2})$ (note the sign change from $R(u) = 1 - \hbar P/u + \cdots$).

For the Yang $R$-matrix $R(u) = 1 - \hbar P/u$, the inverse is $R^{-1}(u) = u/(u - \hbar P) = \sum_{n \geq 0} (\hbar P/u)^n$, where $P^2 = 1$ gives $R^{-1}(u) = 1 + \hbar P/u + \hbar^2/u^2 + \hbar^3 P/u^3 + \cdots$.  The Koszul dual RTT relation is the same quadratic relation with $\hbar \to -\hbar$.
\end{theorem}

\begin{proof}
\textbf{Step 1: Quadratic dual.}
The RTT algebra is quadratic: $V = \mathrm{span}\{T_{ij}^{(r)}\}$, $R = \mathrm{span}\{[\text{RTT relations}]\} \subset V \otimes V$.  The quadratic Koszul dual has generators $V^*$ and relations $R^\perp$ (the orthogonal complement under the pairing $V^* \otimes V^* \times V \otimes V \to \mathbb{C}$).

\textbf{Step 2: Orthogonal complement computation.}
The RTT relation $R_{12}(u-v) T_1 T_2 = T_2 T_1 R_{12}(u-v)$ can be written as:
\[
(1 - \hbar P/(u-v)) T_1 T_2 = T_2 T_1 (1 - \hbar P/(u-v))
\]
Expanding: $T_1 T_2 - T_2 T_1 = \frac{\hbar}{u-v}(P T_1 T_2 - T_2 T_1 P)$.

The orthogonal complement $R^\perp$ consists of quadratic expressions $Q(T^*, T^*)$ that pair to zero with all elements of $R$.  For $\mathfrak{sl}_2$ (where $V = \mathbb{C}^2$ and $P$ is the $4 \times 4$ permutation matrix), the orthogonal complement is computed by direct matrix calculation: the RTT relations span a subspace of $\mathrm{End}(V)^{\otimes 2}$ of codimension equal to the dimension of the space of $R^{-1}$-RTT relations, and $R^\perp$ is exactly the space of $R^{-1}$-RTT relations.  This is verified as follows.  The RTT relation at fixed spectral parameters defines a linear map $\Phi_R\colon \mathrm{End}(V)^{\otimes 2} \to \mathrm{End}(V)^{\otimes 2}$ by $X \mapsto R_{12} X - \sigma(X) R_{12}$ where $\sigma$ is the flip.  Its kernel is the space of solutions, and its image is the space of relations~$R$.  The adjoint map $\Phi_R^*$ with respect to the trace pairing satisfies $\Phi_R^* = \Phi_{R^{-1}}$ (since $R^{-1}$ is the transpose-inverse in the appropriate sense), giving $R^\perp = \ker(\Phi_{R^{-1}}^*) = \mathrm{im}(\Phi_{R^{-1}})$, which is exactly the $R^{-1}$-RTT relations.

\textbf{Step 2a: Local finiteness.}
The Yangian has infinitely many generators $\{T_{ij}^{(r)}\}_{r \geq 0}$ graded by level~$r$, but each graded component $V_r = \mathrm{span}\{T_{ij}^{(r)} : 1 \leq i,j \leq \dim V\}$ is finite-dimensional: $\dim V_r = (\dim V)^2 = 4$ for $\mathfrak{sl}_2$.  The RTT relation~\eqref{eq:rtt-modes} is homogeneous of degree $r + s$ in the level grading (relating generators of levels $r, s$ on both sides), so the quadratic algebra $A = T(V)/\langle R \rangle$ is locally finite: each graded component $A_n = \bigoplus_{r_1 + \cdots + r_k = n} V_{r_1} \cdots V_{r_k} / R_n$ is a quotient of a finite-dimensional space.

\textbf{Step 2b: Graded Koszul duality.}
Since each graded component is finite-dimensional, the Positselski--Polishchuk framework \cite{PP05} for Koszul duality of locally finite graded algebras applies directly.  Their Theorem~2.1 establishes the bar-cobar adjunction and the Koszul property for quadratic algebras that are locally finite in a non-negative grading (which is satisfied here, with the level grading $r \geq 0$).  The key hypothesis --- that the grading is connected ($A_0 = \mathbb{C}$) --- holds after passing to the augmentation ideal.

\textbf{Step 3: $R$-matrix inversion.}
For the Yang $R$-matrix, $R(u) = (u - \hbar P)/u$ (in the defining representation of $\mathfrak{sl}_2$).  The inverse is $R^{-1}(u) = u/(u - \hbar P) = 1 + \hbar P/u + (\hbar P)^2/u^2 + \cdots$.  Since $P^2 = 1$ on $V \otimes V$, this simplifies to $R^{-1}(u) = u/(u-\hbar)$ on the symmetric subspace and $R^{-1}(u) = u/(u+\hbar)$ on the antisymmetric subspace.

\textbf{Step 4: Identification.}
The RTT algebra for $R^{-1}(u)$ is isomorphic to the RTT algebra for $R(u)$ with $\hbar \to -\hbar$.  To see this, note that $R^{-1}(u; \hbar) = 1 + \hbar P/u + O(u^{-2})$ agrees with $R(u; -\hbar) = 1 + \hbar P/u$ at leading order in $1/u$; the higher-order terms $\hbar^{2k}/u^{2k}$ and $\hbar^{2k+1}P/u^{2k+1}$ in $R^{-1}$ do \emph{not} appear in $R(u;-\hbar)$, but they are irrelevant: the mode expansion~\eqref{eq:rtt-modes} extracts only the leading-order data, and the higher-order terms are automatic consequences of the quadratic RTT relation via the Yang--Baxter equation.  Therefore $Y(\mathfrak{sl}_2)^! \cong Y(\mathfrak{sl}_2)^{\hbar \to -\hbar}$, confirming that the Yangian is ``almost self-dual'' up to a sign flip of the deformation parameter.
\end{proof}

\begin{corollary}[Yangian Classical Limit; \ClaimStatusProvedHere]\label{cor:yangian-classical-self-dual}
At $\hbar = 0$ (the classical limit), $R(u) = R^{-1}(u) = 1$, and the RTT algebra becomes commutative: $[T_{ij}^{(r)}, T_{kl}^{(s)}] = 0$, i.e., $\mathrm{Sym}(V)$ where $V = \bigoplus_{r \geq 0} \mathfrak{gl}_2^{(r)}$.  The Koszul dual of the symmetric algebra is the exterior algebra: $\mathrm{Sym}(V)^! = \bigwedge(V^*)$.  Since $R^{-1}(u;0) = R(u;0) = 1$, the orthogonal complement of the defining relations coincides with the relations themselves (both are the full symmetric tensor), confirming the Yangian duality reduces to $\mathrm{Sym} \leftrightarrow \bigwedge$ at $\hbar = 0$.  (In the Drinfeld presentation, the classical limit is the noncommutative algebra $U(\mathfrak{g}[t])$; the two presentations yield different classical limits.)
\end{corollary}

\begin{remark}[Connection to Langlands duality]\label{rem:yangian-langlands}
The sign flip $\hbar \to -\hbar$ in the Yangian Koszul dual is the analog of the Feigin--Frenkel level shift $k \to -k - 2h^\vee$ for affine algebras.  For simply-laced $\mathfrak{g}$, the Yangian $Y(\mathfrak{g})$ is isomorphic to $Y(\mathfrak{g}^\vee)$ (since $\mathfrak{g} \cong \mathfrak{g}^\vee$), and the Koszul duality becomes a form of self-duality.  For non-simply-laced $\mathfrak{g}$, one expects $Y(\mathfrak{g})^! \simeq Y(\mathfrak{g}^\vee)$, realizing Langlands duality at the Yangian level.
\end{remark}

\section{Shifted Yangians}
\label{sec:shifted-yangians}

\subsection{Definition}

\begin{definition}[Shifted Yangian]\label{def:shifted-yangian}
For a coweight $\mu \in X_*(\mathfrak{h})$, the \textbf{shifted Yangian} $Y_\mu(\mathfrak{g})$
(Brundan--Kleshchev, Finkelberg--Tsymbaliuk, KWWY) is generated by elements
$E_i^{(r)}, F_i^{(r)}, H_i^{(r)}$ for $i = 1, \ldots, \mathrm{rank}(\mathfrak{g})$
and $r \geq 0$, satisfying the same Yangian relations as $Y(\mathfrak{g})$ except that the
generating function $H_i(u) = \sum_{r \geq 0} H_i^{(r)} u^{-r}$ is constrained to be a rational
function whose numerator and denominator degrees are prescribed by $\mu$:
\[
H_i(u) = \frac{\prod_{j=1}^{\mu_i^+}(u - a_j^{(i)})}{\prod_{j=1}^{\mu_i^-}(u - b_j^{(i)})}
\]
where $\mu_i^\pm = \langle \alpha_i, \mu^\pm \rangle$ and $\mu = \mu^+ - \mu^-$ decomposes the coweight into dominant parts. The parameters $a_j^{(i)}, b_j^{(i)}$ are complex numbers.
\end{definition}

\begin{conjecture}[Shifted Yangian as $\Eone$-Chiral]\label{conj:shifted-yangian-e1}
The shifted Yangian admits a chiral algebra structure $Y_\mu(\mathfrak{g})^{\mathrm{ch}}$
that is $\Eone$-chiral.  The shift $\mu$ modifies the central terms in the OPE,
and the rational constraint on $H_i(u)$ introduces additional poles in the vertex
operators.
\end{conjecture}


\section{Coulomb Branch Algebras}
\label{sec:coulomb-branch}

\subsection{Definition from Gauge Theory}

\begin{definition}[Coulomb Branch Algebra]\label{def:coulomb-branch}
For a 3d $\mathcal{N} = 4$ gauge theory with gauge group $G$ and matter representation $N$, 
the \textbf{Coulomb branch algebra} $\mathcal{A}(G, N)$ is the quantization of the 
Coulomb branch $\mathcal{M}_C(G, N)$.

For type A quiver gauge theories, $\mathcal{A}(G, N)$ is a shifted Yangian or 
truncated shifted Yangian.
\end{definition}

\begin{theorem}[BFN Construction; \ClaimStatusProvedElsewhere]\label{thm:bfn}
(Braverman--Finkelberg--Nakajima \cite{BFN18}) The Coulomb branch algebra is:
\[
\mathcal{A}(G, N) = H^{G_{\mathcal{O}} \times \mathbb{C}^\times}_*(\mathcal{R}(G, N))
\]
the equivariant Borel--Moore homology of the space of triples $\mathcal{R}(G, N)$.
\end{theorem}

\begin{example}[Quiver of Type $A_n$]\label{ex:coulomb-an}
For the $A_n$ quiver with dimension vector $(1, 2, \ldots, n)$:
\[
\mathcal{A} = Y_{(n-1, n-2, \ldots, 1, 0)}(\mathfrak{sl}_n)
\]
a shifted Yangian.
\end{example}


\section{Cohomological Hall Algebras}
\label{sec:coha}

\subsection{Definition}

\begin{definition}[Cohomological Hall Algebra]\label{def:coha}
For a quiver $Q$ with dimension vector $\mathbf{d}$, the \textbf{cohomological Hall algebra} (CoHA) is:
\[
\mathcal{H}_Q = \bigoplus_{\mathbf{d}} H^{G_{\mathbf{d}}}_{\mathrm{BM},*}(\mathcal{M}_{\mathbf{d}}^{\mathrm{nil}})
\]
where $\mathcal{M}_{\mathbf{d}}^{\mathrm{nil}}$ is the moduli of nilpotent representations.

The product is defined via the Hecke correspondence:
\[
\mathcal{M}_{\mathbf{d}_1} \times \mathcal{M}_{\mathbf{d}_2} \xleftarrow{p} \mathcal{E} \xrightarrow{q} \mathcal{M}_{\mathbf{d}_1 + \mathbf{d}_2}
\]
as $a \star b = q_* p^*(a \boxtimes b)$.
\end{definition}

\begin{conjecture}[CoHA as $\Eone$-Chiral]\label{conj:coha-e1}
The CoHA $\mathcal{H}_Q$ carries an $\Eone$-chiral algebra structure.  The evidence:
\begin{enumerate}[label=(\roman*)]
\item The Hecke product is associative but not commutative.
\item The factorization structure comes from stratifying by dimension vectors.
\item The chiral enhancement should use configuration spaces of points on curves
parameterizing deformation directions.
\end{enumerate}
A rigorous construction requires specifying the vertex operators and verifying
weak associativity (the Borcherds identity).
\end{conjecture}

\begin{example}[Jordan Quiver CoHA]\label{ex:jordan-coha}
For the Jordan quiver (one vertex, one loop), the CoHA is:
\[
\mathcal{H}_{\text{Jordan}} = \bigoplus_{n \geq 0} H^{GL_n}_{\mathrm{BM},*}(\mathcal{N}_n) \;\curvearrowright\; \bigoplus_{n \geq 0} H^*_T(\mathrm{Hilb}^n(\bC^2))
\]
The CoHA acts on the equivariant cohomology of the Hilbert scheme via Nakajima-type operators (Schiffmann--Vasserot), realizing a Heisenberg/W-algebra structure on the target.  The CoHA itself carries a \emph{different} product structure (the shuffle product) from the OPE product on the chiral algebra side.
\end{example}
